\documentclass[11pt,oneside]{article}

\usepackage[a4paper,margin=26mm,headheight=15pt]{geometry}
\usepackage[T1]{fontenc}
\usepackage{lmodern}
\usepackage{microtype}
\usepackage{amsmath,amssymb,amsthm,mathtools,mathrsfs,bm}
\usepackage{booktabs,longtable,tabularx,array,multirow}
\usepackage{graphicx,float}
\usepackage{xcolor}
\usepackage{enumitem}
\usepackage{fancyhdr}
\usepackage{titlesec}
\usepackage{xurl}
\usepackage{hyperref}
\usepackage[nameinlink,noabbrev]{cleveref}

\definecolor{linkblue}{RGB}{25,75,135}
\definecolor{shadegray}{RGB}{246,247,249}
\definecolor{rulegray}{RGB}{190,197,205}
\hypersetup{
  colorlinks=true,
  linkcolor=linkblue,
  citecolor=linkblue,
  urlcolor=linkblue,
  pdftitle={Beyond the Dyadic Fabius Web},
  pdfauthor={Research synthesis for the ProveIt Fabius-function corpus},
  pdfsubject={Negative-q conjugacy, q-power decimation, sinc--Pochhammer factorization, and Thue--Morse reciprocity},
  pdfkeywords={Fabius function, Rvachev up-function, Thue-Morse, q-binomial, q-Pochhammer, Bernoulli, Bell, Lambert W}
}

\pagestyle{fancy}
\fancyhf{}
\fancyhead[L]{\small Beyond the dyadic Fabius web}
\fancyhead[R]{\small\nouppercase{\leftmark}}
\fancyfoot[C]{\thepage}
\renewcommand{\headrulewidth}{0.4pt}
\titleformat{\section}{\Large\bfseries}{\thesection}{0.7em}{}
\titleformat{\subsection}{\large\bfseries}{\thesubsection}{0.7em}{}
\titleformat{\subsubsection}{\normalsize\bfseries}{\thesubsubsection}{0.7em}{}
\setcounter{secnumdepth}{3}
\setcounter{tocdepth}{2}
\setlist[itemize]{leftmargin=2em,itemsep=0.25em,topsep=0.4em}
\setlist[enumerate]{leftmargin=2.3em,itemsep=0.25em,topsep=0.4em}
\setlength{\emergencystretch}{2em}
\allowdisplaybreaks

\newtheorem{theorem}{Theorem}[section]
\newtheorem{proposition}[theorem]{Proposition}
\newtheorem{lemma}[theorem]{Lemma}
\newtheorem{corollary}[theorem]{Corollary}
\newtheorem{conjecture}[theorem]{Conjecture}
\theoremstyle{definition}
\newtheorem{definition}[theorem]{Definition}
\newtheorem{example}[theorem]{Example}
\newtheorem{problem}[theorem]{Research problem}
\theoremstyle{remark}
\newtheorem{remark}[theorem]{Remark}
\newtheorem{warning}[theorem]{Warning}

\newcommand{\R}{\mathbb R}
\newcommand{\C}{\mathbb C}
\newcommand{\N}{\mathbb N}
\newcommand{\Z}{\mathbb Z}
\newcommand{\E}{\mathbb E}
\newcommand{\Prob}{\mathbb P}
\newcommand{\Var}{\operatorname{Var}}
\newcommand{\Law}{\mathcal L}
\newcommand{\Unif}{\operatorname{Unif}}
\newcommand{\supp}{\operatorname{supp}}
\newcommand{\sinc}{\operatorname{sinc}}
\newcommand{\Up}{\operatorname{up}}
\newcommand{\ii}{\mathrm i}
\newcommand{\e}{\mathrm e}
\newcommand{\dd}{\,\mathrm d}
\newcommand{\bigO}{\mathcal O}
\newcommand{\smallo}{o}
\newcommand{\qbinom}[3]{\genfrac{[}{]}{0pt}{}{#1}{#2}_{#3}}
\newcommand{\Bell}{\mathcal B}
\newcommand{\epsTM}{\varepsilon}
\newcommand{\Dsc}{\mathcal D}
\newcommand{\abs}[1]{\left\lvert#1\right\rvert}
\newcommand{\norm}[1]{\left\lVert#1\right\rVert}
\newcommand{\floor}[1]{\left\lfloor#1\right\rfloor}
\newcommand{\ceil}[1]{\left\lceil#1\right\rceil}
\newcommand{\repo}{\texttt{ProveIt}}

\newenvironment{statusbox}{%
  \par\smallskip\noindent\begin{center}\begin{minipage}{0.94\textwidth}%
  \hrule\vspace{0.6em}\small
}{%
  \vspace{0.6em}\hrule\end{minipage}\end{center}\smallskip
}

\newcommand{\Corpus}{\textsc{Corpus}}
\newcommand{\Derived}{\textsc{Derived here}}
\newcommand{\Conjectural}{\textsc{Conjectural}}

\title{\bfseries Beyond the Dyadic Fabius Web\\[0.35em]
\large Negative-\(q\) conjugacy, \(q\)-power decimation, sinc--Pochhammer factorization,\\
Thue--Morse desingularization, and reciprocal-radix limits}
\author{Research synthesis for the \texttt{ProveIt} Fabius-function corpus}
\date{28 August 2026}

\begin{document}
\maketitle

\begin{abstract}
The Fabius function and the Rvachev up-function sit at an unusual meeting point of
smooth compactly supported probability densities, dyadic functional equations,
Thue--Morse cancellation, infinite sinc products, Gaussian \(q\)-binomial
coefficients, \(q\)-Pochhammer symbols, Bernoulli cumulants, Bell polynomials, and
Lambert--\(W\) endpoint asymptotics.  This report audits the current and archived
LaTeX families in the \repo{} Fabius-function repository and then develops a
single geometric-uniform probability family
\[
 X_q=(1-q)\sum_{n\geq0}q^nU_n,\qquad U_n\stackrel{\mathrm{iid}}{\sim}\Unif[0,1],
 \qquad |q|<1,
\]
as a common organizing object.

Several exact results appear to be absent from the inspected corpus.  For
\(0<r<1\), the negative-parameter law is not a genuinely new shape:
\[
 X_{-r}-\tfrac12\ \stackrel d=\ \frac{1+r}{1-r}
 \left(X_r-\tfrac12\right).
\]
Consequently the variance-normalized and support-normalized laws at \(r\) and
\(-r\) coincide.  Splitting digit indices modulo \(m\) gives a second exact
identity,
\[
 X_q\stackrel d=\sum_{j=0}^{m-1}\frac{q^j}{[m]_q}X_{q^m}^{(j)},
\]
so, in particular, the ordinary Fabius law is the convolutional barycenter of
two independent quarter-base laws.  Euler's sine product then yields the
sinc--Pochhammer factorization
\[
 \prod_{n\ge0}\sinc\!\left(\frac{(1-q)q^nt}{2}\right)
 =\prod_{k\ge1}\left(\frac{(1-q)^2t^2}{4\pi^2k^2};q^2\right)_\infty.
\]
This formula packages the zero set, spectral zeta function, Bernoulli
cumulants, and the equality of normalized \(q\) and \(-q\) profiles.

A further reciprocal-radix theorem incorporates the requested exterior values
\(q=\pm2,\pm4\).  Normalized finite sums at \(b\) and \(b^{-1}\) have exactly the
same law after reversing the iid digits.  More strikingly, a desingularized
finite Thue--Morse exponential polynomial is exactly the centered moment
generating function of that finite convolution.  Let
\[
 T_{m,b}(s)=\prod_{j=0}^{m-1}(1-\e^{b^js})
 =\sum_{n<2^m}(-1)^{s_2(n)}\e^{sA_b(n)},
 \qquad A_b(n)=\sum_{j=0}^{m-1}b^jb_j(n),
\]
and \(\alpha_{m,b}=(1-b)/(1-b^m)\).  Then
\[
 \E\e^{t(X_b^{[m]}-1/2)}
 =\frac{(-1)^m\e^{-t/2}T_{m,b}(\alpha_{m,b}t)}
 {\alpha_{m,b}^m b^{\binom m2}t^m}.
\]
For \(|b|>1\), the left side converges to the centered transform of
\(X_{1/b}\).  Thus \(b=2,-2,4,-4\) produce, respectively, the laws at
\(q=1/2,-1/2,1/4,-1/4\).  The report also proves log-concavity, an exact plateau
phase for \(|q|\leq1/2\), reciprocal \(q\)-Lagrange row reversal, a
position-weighted Thue--Morse/\(q\)-Pochhammer master identity, and general-radix
Bell--Bernoulli formulas.  Conjectural sections formulate a general-\(q\)
periodic Lambert--\(W\) transseries and a strict-log-concavity phase beyond the
plateau threshold.  A fully commented Python supplement verifies the exact and
analytic identities and produces the figures.
\end{abstract}

\begin{statusbox}
\textbf{Novelty and proof-status convention.}
\Corpus{} means that a result or its essential mechanism is already present in
the inspected repository or cited literature.  \Derived{} means that the
statement was derived in this audit and was not located by exact-term,
structural, and formula searches across the current canonical documents,
consolidated frontier volumes, draft syntheses, and archive manifests.  This is
a claim of novelty relative to the inspected corpus, not a claim of global
priority.  \Conjectural{} marks statements for which only heuristics, numerical
evidence, or a proof program is supplied.  None of the newly derived claims is
presented as Lean-formalized or peer-reviewed.
\end{statusbox}

\tableofcontents
\clearpage

\section{Scope, corpus audit, and the non-duplication boundary}

\subsection{What was inspected}
The recursive audit began with every current LaTeX-bearing family under
\url{https://github.com/VladimirReshetnikov/ProveIt/tree/main/Analysis/FabiusFunction/docs}.
The principal sources were the formal exposition
\cite{repo-primary}, the consolidated non-formalized frontier volume
\cite{repo-frontiers}, the exponents-and-\(q\)-series synthesis
\cite{repo-qseries}, the Thue--Morse atlas and frontier synthesis
\cite{repo-thuemorse}, and the archive map \cite{repo-archive}.  The archive
manifest was used to traverse historical families while treating byte-identical
or explicitly consolidated copies as provenance rather than as independent
mathematics.  The repository's paper collection and its surrounding citations
were cross-checked against Arias de Reyna's analytic and arithmetic papers
\cite{arias-up,arias-arithmetic}, Haugland's dyadic evaluation paper
\cite{haugland}, the digital and \(q\)-digital binomial theorems
\cite{nguyen-digital,mansour-nguyen}, and later Thue--Morse work
\cite{aistleitner-hofer-larcher,toth}.

The current primary exposition already contains a broad weighted-uniform-series
API, the geometric law for \(|q|<1\), its affine recurrence, reflection,
absolute continuity, and positive-\(q\) density bridges to Fabius and Rvachev.
It also contains generic Gaussian-binomial and geometric Lagrange machinery,
including explicit quarter-base cases.  The current frontier register itself
identifies negative-parameter density interpretation, the general-\(q\)
Fourier route, centered transport, shape, moments, transforms, derivative
hierarchies, and endpoint flatness as remaining boundaries
\cite{repo-frontiers}.  The \(q\)-series synthesis develops general
exponent-sequence sinc products, zero multiplicities, spectral zeta functions,
cumulants, and \(q\)-Richardson ideas, while the Thue--Morse synthesis develops
Boolean-cube products, Bell--Bernoulli moment formulas, Gaussian-binomial
interpolation, and weighted digit enumerators \cite{repo-qseries,repo-thuemorse}.

\subsection{Results deliberately not counted as new}
The following are used as foundations rather than advertised as discoveries:
\begin{itemize}
\item the identification of the half-base geometric-uniform law with the
  Fabius CDF and the Rvachev density;
\item the ordinary Fourier product for the up-function and its finite-product
  approximants;
\item standard \(q\)-Pochhammer and Gaussian-binomial identities, including
  the finite \(q\)-binomial theorem;
\item generic cumulant/Bell-polynomial conversion and generic spectral-zeta
  formulas for exponent-sequence sinc products;
\item the classical Thue--Morse product
  \(\prod_{j<m}(1-z^{2^j})=\sum_{n<2^m}(-1)^{s_2(n)}z^n\);
\item the fixed-Hamming-weight Gaussian-binomial enumerator, which is closely
  related to the \(q\)-digital binomial theorem of Mansour and Nguyen
  \cite{mansour-nguyen}.
\end{itemize}

\subsection{Principal new-to-corpus seams}
The report concentrates on interactions not found in the audit:
\begin{enumerate}
\item affine sign conjugacy for \(q\mapsto-q\), including exact density, CDF,
  quantile, support, and inverse-function transport;
\item residue-class decimation \(q\mapsto q^m\), especially the exact
  \(1/2\leftrightarrow1/4\) convolution identity;
\item the double infinite sinc--\(q\)-Pochhammer product and its compatibility
  with decimation;
\item exact finite reciprocal-radix duality for \(q\leftrightarrow q^{-1}\),
  giving a probability interpretation to \(q=\pm2,\pm4\);
\item a desingularized Thue--Morse transform representation of finite and
  limiting geometric-uniform laws;
\item log-concavity and the exact plateau phase for the whole real family;
\item a two-variable digit-position master product whose negative-\(q\)
  specialization replaces Hamming parity by position parity;
\item general-radix Bell--Bernoulli formulas and centered \(b\mapsto-b\)
  conjugacy;
\item a general-\(q\) Lambert--\(W\) endpoint program with exact negative-
  parameter transport.
\end{enumerate}

\section{The geometric-uniform family as a common language}

\subsection{Definition and elementary structure}
Let \((U_n)_{n\ge0}\) be independent \(\Unif[0,1]\) random variables.  For real
\(|q|<1\), define
\begin{equation}\label{eq:Xq-def}
 X_q=(1-q)\sum_{n=0}^{\infty}q^nU_n.
\end{equation}
The series converges absolutely almost surely.  Its coefficients sum to one,
so \(\E X_q=1/2\).  Writing \(V_n=U_n-1/2\), one obtains the centered form
\begin{equation}\label{eq:centered-series}
 Y_q:=X_q-\frac12=(1-q)\sum_{n=0}^{\infty}q^nV_n,
 \qquad V_n\stackrel{\mathrm{iid}}{\sim}\Unif[-1/2,1/2].
\end{equation}
The symmetry of each \(V_n\) gives \(Y_q\stackrel d=-Y_q\), hence the law of
\(X_q\) is symmetric about \(1/2\) for every real \(|q|<1\), including negative
\(q\).

Separating the first digit gives the affine distributional fixed point
\begin{equation}\label{eq:affine-fixed-point}
 X_q\stackrel d=(1-q)U+qX_q',
\end{equation}
where \(U\sim\Unif[0,1]\), \(X_q'\stackrel d=X_q\), and the two variables are
independent.

\begin{proposition}[Wasserstein uniqueness; \Derived]\label{prop:wasserstein}
For \(|q|<1\), the map
\[
 \mathsf T_q(\nu)=\Law((1-q)U+qY),\qquad Y\sim\nu,
\]
is a contraction on probability measures with finite first moment in the
\(1\)-Wasserstein metric:
\[
 W_1(\mathsf T_q\mu,\mathsf T_q\nu)\le |q|W_1(\mu,\nu).
\]
Consequently \(\Law(X_q)\) is the unique finite-first-moment solution of
\eqref{eq:affine-fixed-point}.
\end{proposition}

\begin{proof}
Choose a coupling \((Y,Z)\) of \(\mu,\nu\), independent of the same uniform
variable \(U\).  Then
\[
 \E\abs{(1-q)U+qY-((1-q)U+qZ)}=|q|\E\abs{Y-Z}.
\]
Taking the infimum over couplings gives the contraction estimate.  Banach's
fixed-point principle gives uniqueness.  The partial sums in
\eqref{eq:Xq-def} converge to the fixed point, proving existence.
\end{proof}

\subsection{Support}
For positive \(q=r\), all coefficients in \eqref{eq:Xq-def} are nonnegative
and sum to one, so \(\supp X_r=[0,1]\).  For \(q=-r\), positive coefficients
occur at even indices and negative coefficients at odd indices.  Their sums are
\[
 \sum_{j\ge0}(1+r)r^{2j}=\frac1{1-r},\qquad
 -\sum_{j\ge0}(1+r)r^{2j+1}=-\frac r{1-r}.
\]
Thus
\begin{equation}\label{eq:negative-support}
 \supp X_{-r}=\left[-\frac r{1-r},\frac1{1-r}\right].
\end{equation}
The midpoint is still \(1/2\).

\subsection{Fabius and Rvachev normalization}
Let \(F_q\) and \(f_q\) denote the CDF and density when these are defined.  At
\(q=1/2\), the repository's normalization gives
\begin{equation}\label{eq:fabius-bridge}
 F_{1/2}(x)=F(x),\qquad
 f_{1/2}(x)=2\Up(2x-1),
\end{equation}
where \(F\) is the bounded Fabius function and \(\Up\) is the standard Rvachev
bump supported on \([-1,1]\) \cite{repo-primary,arias-up}.

For general \(q\ne0\), it is convenient to use the support half-width
\begin{equation}\label{eq:Rq}
 R_q=\frac{1-q}{2(1-|q|)}
\end{equation}
and the support-normalized density
\begin{equation}\label{eq:Uq-def}
 \mathcal U_q(x)=R_q f_q\!\left(\frac12+R_qx\right),\qquad -1\le x\le1.
\end{equation}
At \(q=1/2\), \(R_q=1/2\), and \eqref{eq:fabius-bridge} gives
\(\mathcal U_{1/2}=\Up\).

\section{Negative \texorpdfstring{\(q\)}{q}: exact affine sign conjugacy}

\subsection{The sign-conjugacy theorem}
\begin{theorem}[Affine sign conjugacy; \Derived]\label{thm:sign-conjugacy}
Let \(0<r<1\) and
\[
 c_r=\frac{1+r}{1-r}.
\]
Then
\begin{equation}\label{eq:sign-law}
 X_{-r}-\frac12\ \stackrel d=\ c_r\left(X_r-\frac12\right).
\end{equation}
Equivalently, for the centered characteristic functions
\(\psi_q(t)=\E\e^{\ii t(X_q-1/2)}\),
\begin{equation}\label{eq:sign-char}
 \psi_{-r}(t)=\psi_r(c_rt).
\end{equation}
\end{theorem}

\begin{proof}
From \eqref{eq:centered-series},
\[
 Y_{-r}=(1+r)\sum_{n\ge0}(-r)^nV_n.
\]
Set \(W_n=(-1)^nV_n\).  Since the centered uniform law is symmetric, the
variables \(W_n\) are again independent \(\Unif[-1/2,1/2]\).  Therefore
\[
 Y_{-r}\stackrel d=(1+r)\sum_{n\ge0}r^nW_n
 =\frac{1+r}{1-r}\left((1-r)\sum_{n\ge0}r^nW_n\right)
 \stackrel d=c_rY_r.
\]
Taking characteristic functions gives \eqref{eq:sign-char}.
\end{proof}

This elementary sign change closes several negative-parameter gaps at once.
It also explains why all even normalized invariants depend on \(q^2\).

\begin{corollary}[CDF, density, and quantile transport; \Derived]
For \(0<r<1\),
\begin{align}
 F_{-r}(x)
 &=F_r\!\left(\frac12+\frac{x-1/2}{c_r}\right),\label{eq:cdf-transport}\\
 f_{-r}(x)
 &=\frac1{c_r}f_r\!\left(\frac12+\frac{x-1/2}{c_r}\right),\label{eq:density-transport}\\
 F_{-r}^{-1}(u)
 &=\frac12+c_r\left(F_r^{-1}(u)-\frac12\right),\qquad 0<u<1.
 \label{eq:quantile-transport}
\end{align}
Every interior or endpoint asymptotic expansion for \(F_r\), \(f_r\), or the
inverse CDF transfers exactly to \(-r\) by these affine substitutions.
\end{corollary}

\begin{proof}
Equation \eqref{eq:sign-law} is an increasing affine transformation because
\(c_r>0\).  The CDF, density, and generalized inverse formulas follow by the
ordinary change-of-variables rule.
\end{proof}

\begin{example}[The values \(q=-1/2\) and \(q=-1/4\)]
At \(r=1/2\), \(c_r=3\), so
\begin{equation}\label{eq:minus-half-up}
 F_{-1/2}(x)=F\!\left(\frac{x+1}{3}\right),\qquad
 f_{-1/2}(x)=\frac23\Up\!\left(\frac{2x-1}{3}\right),
\end{equation}
with support \([-1,2]\).  At \(r=1/4\), \(c_r=5/3\), so
\begin{equation}\label{eq:minus-quarter}
 F_{-1/4}(x)=F_{1/4}\!\left(\frac{3x+1}{5}\right),\qquad
 f_{-1/4}(x)=\frac35f_{1/4}\!\left(\frac{3x+1}{5}\right),
\end{equation}
with support \([-1/3,4/3]\).
\end{example}

\subsection{Normalized-shape collapse}
The variance follows immediately from independent summation:
\begin{equation}\label{eq:variance}
 \sigma_q^2:=\Var(X_q)
 =\frac1{12}(1-q)^2\sum_{n\ge0}q^{2n}
 =\frac{1-q}{12(1+q)}.
\end{equation}

\begin{corollary}[Variance-normalized collapse; \Derived]\label{cor:standard-collapse}
Let \(Z_q=(X_q-1/2)/\sigma_q\).  For \(0<r<1\),
\begin{equation}\label{eq:standard-collapse}
 Z_{-r}\stackrel d=Z_r.
\end{equation}
Hence every standardized moment, cumulant, quantile, entropy shift, and shape
statistic is an even function of \(q\).
\end{corollary}

\begin{proof}
Equation \eqref{eq:variance} gives \(\sigma_{-r}=c_r\sigma_r\).  Divide
\eqref{eq:sign-law} by this identity.
\end{proof}

\begin{corollary}[Support-normalized Rvachev profile; \Derived]
For \(0<r<1\),
\begin{equation}\label{eq:profile-collapse}
 \mathcal U_{-r}=\mathcal U_r.
\end{equation}
Thus changing the sign of \(q\) changes only the physical location and scale;
it does not create a new normalized bump.
\end{corollary}

\begin{proof}
For positive \(r\), \(R_r=1/2\).  For negative \(-r\),
\(R_{-r}=c_r/2\).  Substitute \eqref{eq:density-transport} into
\eqref{eq:Uq-def}.
\end{proof}

\section{Residue-class decimation and the \texorpdfstring{\(1/2\leftrightarrow1/4\)}{1/2--1/4} bridge}

\subsection{The general decimation law}
Write the \(q\)-integer
\[
 [m]_q=1+q+\cdots+q^{m-1}=\frac{1-q^m}{1-q}.
\]

\begin{theorem}[\(q\)-power decimation; \Derived]\label{thm:decimation}
Let \(|q|<1\) and \(m\ge1\).  There exist independent copies
\(X_{q^m}^{(0)},\ldots,X_{q^m}^{(m-1)}\) such that
\begin{equation}\label{eq:decimation}
 X_q\stackrel d=\sum_{j=0}^{m-1}a_{m,j}(q)X_{q^m}^{(j)},
 \qquad
 a_{m,j}(q)=\frac{q^j}{[m]_q}
 =\frac{(1-q)q^j}{1-q^m}.
\end{equation}
The coefficients satisfy \(\sum_ja_{m,j}(q)=1\), although some are negative
when \(q<0\).  Consequently
\begin{equation}\label{eq:decimation-char}
 \phi_q(t)=\prod_{j=0}^{m-1}\phi_{q^m}(a_{m,j}(q)t),
 \qquad
 \psi_q(t)=\prod_{j=0}^{m-1}\psi_{q^m}(a_{m,j}(q)t).
\end{equation}
\end{theorem}

\begin{proof}
Split the series \eqref{eq:Xq-def} by residue classes modulo \(m\):
\begin{align*}
 X_q
 &=(1-q)\sum_{j=0}^{m-1}\sum_{\ell\ge0}q^{m\ell+j}U_{m\ell+j}\\
 &=\sum_{j=0}^{m-1}\frac{(1-q)q^j}{1-q^m}
 \left((1-q^m)\sum_{\ell\ge0}(q^m)^\ell U_{m\ell+j}\right).
\end{align*}
The parenthesized variables are independent copies of \(X_{q^m}\).  Their
coefficients sum to one.  Independence gives the characteristic products.
\end{proof}

\begin{corollary}[Cumulant compatibility]
For every \(n\ge1\) for which the cumulants are written,
\begin{equation}\label{eq:decimation-cumulants}
 \kappa_n(q)=\kappa_n(q^m)\sum_{j=0}^{m-1}a_{m,j}(q)^n.
\end{equation}
After inserting the Bernoulli formula from \cref{sec:cumulants}, this becomes
the \(q\)-integer identity
\begin{equation}\label{eq:qinteger-decimation}
 \sum_{j=0}^{m-1}\left(\frac{q^j}{[m]_q}\right)^n
 =\frac{[m]_{q^n}}{[m]_q^n}.
\end{equation}
\end{corollary}

\subsection{Half base as a convolution of quarter-base laws}
For \(q=1/2\), \(m=2\), the coefficients are \(2/3\) and \(1/3\):
\begin{equation}\label{eq:half-quarter}
 X_{1/2}\stackrel d=\frac23X_{1/4}^{(0)}+\frac13X_{1/4}^{(1)}.
\end{equation}
For \(q=-1/2\), the same split gives
\begin{equation}\label{eq:minus-half-quarter}
 X_{-1/2}\stackrel d=2X_{1/4}^{(0)}-X_{1/4}^{(1)}.
\end{equation}
The first identity is a direct probabilistic bridge from the Fabius/Rvachev
law to the quarter-base family already prominent in the repository.  The
second shows that the negative half-base law can be built from the same
quarter-base blocks, with one reflected coefficient.

For a density \(g\) on a centered coordinate, define
\[
 (\Dsc_a g)(x)=\frac1{|a|}g(x/a).
\]
Because the support-normalized quarter-base profile is even, the sign in
\eqref{eq:minus-half-quarter} disappears after normalization.

\begin{corollary}[Rvachev convolutional barycenter; \Derived]\label{cor:up-quarter}
The standard Rvachev bump satisfies
\begin{equation}\label{eq:up-quarter-convolution}
 \Up=\Dsc_{2/3}\mathcal U_{1/4}*\Dsc_{1/3}\mathcal U_{1/4}.
\end{equation}
The same right side is the support-normalized density of \(X_{-1/2}\).
\end{corollary}

\begin{proof}
For positive \(q\), the support radius is \(1/2\), so support normalization
turns \eqref{eq:half-quarter} into
\(2(X_{1/2}-1/2)=\frac23\,2(X_{1/4}^{(0)}-1/2)
+\frac13\,2(X_{1/4}^{(1)}-1/2)\).  Convolution gives
\eqref{eq:up-quarter-convolution}.  For \(-1/2\), divide the centered version
of \eqref{eq:minus-half-quarter} by its support radius \(3/2\); the second
coefficient becomes \(-1/3\), which has the same effect as \(+1/3\) on an even
input density.
\end{proof}

More generally, for every \(m\ge1\),
\begin{equation}\label{eq:half-to-2m}
 X_{1/2}\stackrel d=\frac{1}{2^m-1}\sum_{j=0}^{m-1}2^{m-1-j}
 X_{2^{-m}}^{(j)}.
\end{equation}
At \(m=3\), the weights are \(4/7,2/7,1/7\) and the building block is
\(X_{1/8}\).

\section{Fourier products, smoothness, and endpoint flatness}

\subsection{The centered sinc product}
For \(V\sim\Unif[-1/2,1/2]\),
\(\E\e^{\ii tV}=\sinc(t/2)\), where \(\sinc z=\sin z/z\).  Hence
\begin{equation}\label{eq:psi-sinc}
 \psi_q(t)=\prod_{n=0}^{\infty}
 \sinc\!\left(\frac{(1-q)q^nt}{2}\right),
 \qquad |q|<1.
\end{equation}
The uncentered characteristic function is
\(\phi_q(t)=\e^{\ii t/2}\psi_q(t)\), and the fixed-point equation is
\begin{equation}\label{eq:psi-refinement}
 \psi_q(t)=\sinc\!\left(\frac{(1-q)t}{2}\right)\psi_q(qt).
\end{equation}

\subsection{Superpolynomial Fourier decay}
\begin{theorem}[General real-\(q\) Fourier decay; \Derived]\label{thm:fourier-decay}
Let \(0<|q|<1\), put \(r=|q|\) and \(\lambda=-\log r\).  Then, as
\(|t|\to\infty\),
\begin{equation}\label{eq:fourier-decay}
 \log|\psi_q(t)|
 \le -\frac{(\log|t|)^2}{2\lambda}+\bigO_q(\log|t|).
\end{equation}
In particular, \(|t|^N\psi_q(t)\to0\) for every \(N\ge0\).
\end{theorem}

\begin{proof}
Set \(A=(1-q)/2>0\).  Since
\(|\sinc x|\le\min(1,|x|^{-1})\), choose
\(M=\floor{\log(A|t|)/\lambda}\) when \(A|t|>1\).  The first \(M\) factors in
\eqref{eq:psi-sinc} give
\begin{align*}
 \log|\psi_q(t)|
 &\le -\sum_{n=0}^{M-1}\log(A|t|r^n)\\
 &=-M\log(A|t|)+\frac{\lambda M(M-1)}2.
\end{align*}
Substituting \(M=\log(A|t|)/\lambda+\bigO(1)\) yields
\eqref{eq:fourier-decay}.
\end{proof}

\begin{corollary}[Smooth density and flat endpoints; \Derived]
For every real \(q\) with \(0<|q|<1\), the law of \(X_q\) has a compactly
supported \(C^\infty\) density.  Every derivative is obtained by Fourier
inversion,
\begin{equation}\label{eq:density-derivative}
 f_q^{(k)}(x)=\frac1{2\pi}\int_{\R}(-\ii t)^k\e^{-\ii tx}\phi_q(t)\dd t,
\end{equation}
and all derivatives vanish at both endpoints of the support.
\end{corollary}

\begin{proof}
The decay in \cref{thm:fourier-decay} makes every polynomial multiple of
\(\phi_q\) integrable, so repeated Fourier inversion gives a smooth density.
A smooth function that is identically zero on an exterior open interval has
all one-sided derivatives zero at the support endpoint; smoothness across that
point makes every endpoint derivative vanish.
\end{proof}

The case \(q=0\) is the uniform law and is intentionally excluded: its density
is not smooth at the endpoints.  This isolates a sharp regularity transition at
\(q=0\).

\section{The sinc--\texorpdfstring{\(q\)}{q}-Pochhammer double product}

\subsection{Factorization theorem}
Use the standard notation
\[
 (a;Q)_\infty=\prod_{n=0}^{\infty}(1-aQ^n)
\]
for the \(q\)-Pochhammer symbol \cite{dlmf}.  Euler's product
\(\sinc z=\prod_{k\ge1}(1-z^2/(\pi^2k^2))\) can be inserted into every factor
of \eqref{eq:psi-sinc}.

\begin{theorem}[Sinc--Pochhammer factorization; \Derived]\label{thm:sinc-poch}
For real \(|q|<1\),
\begin{equation}\label{eq:sinc-poch}
 \boxed{
 \psi_q(t)=\prod_{k=1}^{\infty}
 \left(\frac{(1-q)^2t^2}{4\pi^2k^2};q^2\right)_\infty.}
\end{equation}
The product converges locally uniformly in \(t\in\C\), so it represents the
same entire function as the sinc product.
\end{theorem}

\begin{proof}
Euler's product gives
\[
 \sinc\!\left(\frac{(1-q)q^nt}{2}\right)
 =\prod_{k\ge1}\left(1-
 \frac{(1-q)^2q^{2n}t^2}{4\pi^2k^2}\right).
\]
On every compact set in \(t\), the double series of absolute first-order
perturbations is bounded by a constant times
\(\sum_n|q|^{2n}\sum_kk^{-2}<\infty\).  The products may therefore be
reordered, and the product over \(n\) is exactly a \(q^2\)-Pochhammer symbol.
\end{proof}

\begin{example}[The four interior values requested in the prompt]
The factorization specializes to
\begin{align}
 \psi_{1/2}(t)
 &=\prod_{k\ge1}\left(\frac{t^2}{16\pi^2k^2};\frac14\right)_\infty,
 \label{eq:poch-half}\\
 \psi_{-1/2}(t)
 &=\prod_{k\ge1}\left(\frac{9t^2}{16\pi^2k^2};\frac14\right)_\infty
 =\psi_{1/2}(3t),\label{eq:poch-minus-half}\\
 \psi_{1/4}(t)
 &=\prod_{k\ge1}\left(\frac{9t^2}{64\pi^2k^2};\frac1{16}\right)_\infty,
 \label{eq:poch-quarter}\\
 \psi_{-1/4}(t)
 &=\prod_{k\ge1}\left(\frac{25t^2}{64\pi^2k^2};\frac1{16}\right)_\infty
 =\psi_{1/4}\!\left(\frac53t\right).
 \label{eq:poch-minus-quarter}
\end{align}
Thus \(q=\pm1/2\) share Pochhammer base \(1/4\), while \(q=\pm1/4\) share
base \(1/16\).  The sign changes only the spectral scale.
\end{example}

\subsection{Decimation is Pochhammer dissection}
The elementary dissection identity
\begin{equation}\label{eq:poch-dissection}
 (a;Q)_\infty=\prod_{j=0}^{m-1}(aQ^j;Q^m)_\infty
\end{equation}
is the product-theoretic shadow of \cref{thm:decimation}.
Indeed, with
\(A_k=(1-q)^2t^2/(4\pi^2k^2)\),
\begin{align*}
 \psi_{q^m}(a_{m,j}(q)t)
 &=\prod_{k\ge1}
 \left(
 \frac{(1-q^m)^2a_{m,j}(q)^2t^2}{4\pi^2k^2};q^{2m}
 \right)_\infty\\
 &=\prod_{k\ge1}(A_kq^{2j};q^{2m})_\infty.
\end{align*}
Multiplying over \(j\) and using \eqref{eq:poch-dissection} gives exactly
\eqref{eq:sinc-poch}.  Therefore residue-class splitting of the random digits
and residue-class splitting of the Pochhammer product are not analogous
operations; they are the same identity in probability and product coordinates.

\section{Spectral zeta, Bernoulli cumulants, and Bell moments}\label{sec:cumulants}

\subsection{The indexed zero spectrum}
The positive zeros of the individual sinc factors occur at
\begin{equation}\label{eq:zeros}
 \tau_{n,k}=\frac{2\pi k}{(1-q)|q|^n},
 \qquad n\ge0,\quad k\ge1.
\end{equation}
Coincident zeros are counted with their factor multiplicity.  Define the
indexed spectral zeta function
\[
 \mathcal Z_q(s)=\sum_{n\ge0}\sum_{k\ge1}\tau_{n,k}^{-s}.
\]

\begin{proposition}[Geometric spectral zeta; synthesis of corpus machinery]
For \(\Re s>1\),
\begin{equation}\label{eq:spectral-zeta}
 \mathcal Z_q(s)=
 \left(\frac{1-q}{2\pi}\right)^s\frac{\zeta(s)}{1-|q|^s}.
\end{equation}
Near the origin,
\begin{equation}\label{eq:log-zeta}
 \log\psi_q(t)=-\sum_{m\ge1}\frac{\mathcal Z_q(2m)}m t^{2m}.
\end{equation}
\end{proposition}

\begin{proof}
Sum the geometric progression over \(n\) and the Riemann zeta series over
\(k\).  The logarithm of the canonical even product is
\(-\sum_{m\ge1}t^{2m}\sum_{n,k}\tau_{n,k}^{-2m}/m\), proving
\eqref{eq:log-zeta} where the power series converges.
\end{proof}

\subsection{Bernoulli cumulants}
The cumulant generating function of a uniform variable is
\begin{equation}\label{eq:uniform-cgf}
 \log\frac{\e^t-1}{t}
 =\frac t2+\sum_{n=2}^{\infty}\frac{B_n}{n}\frac{t^n}{n!},
\end{equation}
with \(B_1=-1/2\).  Independent weighted summation gives the following closed
formula.

\begin{theorem}[Bernoulli cumulants; corpus specialization completed for real
\(q\)]\label{thm:cumulants}
For \(|q|<1\),
\begin{equation}\label{eq:cumulants}
 \kappa_1(X_q)=\frac12,
 \qquad
 \kappa_n(X_q)=\frac{B_n}{n}\frac{(1-q)^n}{1-q^n},\quad n\ge2.
\end{equation}
All odd cumulants above the mean vanish.  For \(m\ge1\), the even cumulants are
also
\begin{equation}\label{eq:cumulant-zeta}
 \kappa_{2m}(X_q)
 =(-1)^{m+1}2(2m-1)!\,\mathcal Z_q(2m).
\end{equation}
\end{theorem}

\begin{proof}
The cumulant of order \(n\) of a weighted independent sum is the sum of the
coefficient \(n\)-th powers times the input cumulant.  Equation
\eqref{eq:uniform-cgf} gives \(\kappa_n(U)=B_n/n\) for \(n\ge2\), and
\[
 \sum_{j\ge0}((1-q)q^j)^n=\frac{(1-q)^n}{1-q^n}.
\]
Comparing the coefficient of \(t^{2m}\) in \eqref{eq:log-zeta} with
\(\log\psi_q(t)=\sum_{m\ge1}(-1)^m\kappa_{2m}t^{2m}/(2m)!\) gives
\eqref{eq:cumulant-zeta}.
\end{proof}

Equation \eqref{eq:cumulants} may also be written
\[
 \kappa_n(X_q)=\frac{B_n}{n}\frac{(1-q)^{n-1}}{[n]_q}.
\]
This is a direct Bernoulli--Gaussian-\(q\)-integer bridge.

\subsection{Standardized cumulants and the Gaussian edge}
From \eqref{eq:variance} and \eqref{eq:cumulants},
\begin{equation}\label{eq:standard-cumulants}
 \kappa_{2m}(Z_q)=
 \frac{B_{2m}}{2m}12^m\frac{(1-q^2)^m}{1-q^{2m}}.
\end{equation}
This formula makes \cref{cor:standard-collapse} algebraically transparent.
The first two non-Gaussian standardized cumulants are
\begin{align}
 \kappa_4(Z_q)&=-\frac65\frac{1-q^2}{1+q^2},\label{eq:kurtosis}\\
 \kappa_6(Z_q)&=\frac{48}{7}
 \frac{(1-q^2)^2}{1+q^2+q^4}.
 \label{eq:sixth}
\end{align}

\begin{proposition}[Gaussian limit; \Derived]
As \(|q|\to1^-\),
\begin{equation}\label{eq:gaussian-limit}
 Z_q\Longrightarrow N(0,1).
\end{equation}
As \(q\to0\), \(Z_q\) converges to the variance-one uniform distribution on
\([-\sqrt3,\sqrt3]\).
\end{proposition}

\begin{proof}
By sign conjugacy it suffices to take \(q=r\to1^-\).  The largest centered
summand has magnitude at most \((1-r)/2\), while the standard deviation is
asymptotic to \(\sqrt{(1-r)/24}\).  Their ratio tends to zero, so the bounded
triangular array satisfies Lindeberg's condition.  The second assertion follows
directly from \(X_q\to U_0\) in \(L^2\).
\end{proof}

\subsection{Bell-polynomial moments}
Let \(\Bell_n\) denote the complete exponential Bell polynomial, normalized by
\[
 \exp\left(\sum_{j\ge1}x_j\frac{t^j}{j!}\right)
 =\sum_{n\ge0}\Bell_n(x_1,\ldots,x_n)\frac{t^n}{n!}.
\]
Then
\begin{equation}\label{eq:bell-raw}
 \E X_q^n=\Bell_n(\kappa_1(q),\ldots,\kappa_n(q)),
\end{equation}
and the centered moments are obtained by setting \(\kappa_1=0\).  Thus every
moment is an explicit rational expression in Bernoulli numbers and the finite
set of \(q\)-integers \([2]_q,\ldots,[n]_q\).  For rational \(q\), all moments
are rational.

\section{Shape theory: log-concavity and the plateau transition}

\subsection{Log-concavity for every real parameter}
\begin{theorem}[Log-concavity; \Derived]\label{thm:logconcave}
For every real \(q\) with \(0<|q|<1\), the density \(f_q\) is log-concave on its
support.  It is symmetric about \(1/2\), positive on the support interior, and
unimodal.
\end{theorem}

\begin{proof}
Every finite partial sum in \eqref{eq:centered-series} is a convolution of
uniform densities on intervals, possibly reflected.  Uniform interval densities
are log-concave, and convolution preserves log-concavity by the
Pr\'ekopa--Leindler theorem.  The partial-sum laws converge weakly to
\(\Law(Y_q)\); the class of log-concave probability measures is closed under
weak convergence.  By \cref{thm:fourier-decay}, the limit has a smooth density.
A nonzero log-concave density is positive on the interior of its interval
support.  Symmetry was proved from the centered series, and a symmetric
log-concave density is unimodal at the center.
\end{proof}

\subsection{Exact plateaus for \texorpdfstring{\(|q|\le1/2\)}{|q|<=1/2}}
For \(q=r>0\), the fixed point \eqref{eq:affine-fixed-point} is a convolution
of \(rX_r'\) with a uniform interval of length \(1-r\).  Therefore
\begin{equation}\label{eq:density-cdf-positive}
 f_r(x)=\frac1{1-r}
 \left[
 F_r\!\left(\frac xr\right)
 -F_r\!\left(\frac{x-(1-r)}r\right)
 \right],
\end{equation}
where the CDF is extended by \(0\) and \(1\) outside \([0,1]\).

\begin{theorem}[Plateau phase; \Derived]\label{thm:plateau}
Let \(0<r\le1/2\).
\begin{enumerate}
\item For positive parameter,
\begin{equation}\label{eq:positive-plateau}
 f_r(x)=\frac1{1-r}\quad\text{for }x\in[r,1-r],
 \qquad
 \norm{f_r}_\infty=\frac1{1-r}.
\end{equation}
\item For negative parameter,
\begin{equation}\label{eq:negative-plateau}
 f_{-r}(x)=\frac1{1+r}
 \quad\text{for }
 x\in\left[\frac{r^2}{1-r},1-\frac{r^2}{1-r}\right],
 \qquad
 \norm{f_{-r}}_\infty=\frac1{1+r}.
\end{equation}
\end{enumerate}
When \(r<1/2\), the displayed intervals have positive length.  At
\(r=1/2\), each plateau collapses to the single central maximizer.
\end{theorem}

\begin{proof}
If \(x\in[r,1-r]\), then \(x/r\ge1\) and
\((x-(1-r))/r\le0\).  Equation \eqref{eq:density-cdf-positive} therefore equals
\(1/(1-r)\).  Convolution with a probability density and a uniform kernel gives
the matching global upper bound.  The negative statement is the affine image
under \eqref{eq:density-transport}; the interval endpoints simplify to
\(r^2/(1-r)\) and its reflection, and the height is divided by \(c_r\).
\end{proof}

The exact plateau at \(q=1/4\), of height \(4/3\) on \([1/4,3/4]\), and its
negative counterpart, of height \(4/5\) on \([1/12,11/12]\), are visible in
\cref{fig:physical-densities}.

\begin{conjecture}[Strict log-concavity beyond the threshold; \Conjectural]
For \(1/2<r<1\), the support-normalized density \(\mathcal U_r\) is strictly
log-concave on \((-1,1)\).  Equivalently, the only obstruction to strictness in
the real geometric-uniform family is the exact plateau phase
\(0<|q|<1/2\) and its critical collapse at \(|q|=1/2\).
\end{conjecture}

The conjecture is consistent with the numerical profiles and with the Gaussian
limit, but a proof appears to require a strict form of convolution
log-concavity adapted to infinitely many nested uniform scales.

\section{Reciprocal \texorpdfstring{\(q\)}{q}: bringing \texorpdfstring{\(\pm2,\pm4\)}{+/-2,+/-4} into the picture}

\subsection{Finite Pochhammer and Gaussian reciprocity}
For \(q\ne0\), finite products satisfy
\begin{equation}\label{eq:poch-reciprocity}
 (a;q)_n=(-a)^nq^{\binom n2}(a^{-1};q^{-1})_n.
\end{equation}
The Gaussian polynomial consequently obeys
\begin{equation}\label{eq:qbinom-reciprocity}
 \qbinom{n}{k}{q}=q^{k(n-k)}\qbinom{n}{k}{q^{-1}}.
\end{equation}
For the nodal polynomial
\(P_n^{(q)}(z)=\prod_{j=1}^n(z-q^j)\),
\begin{equation}\label{eq:nodal-reciprocity}
 P_n^{(q)}(z)=(-1)^nq^{n(n+1)/2}z^nP_n^{(q^{-1})}(z^{-1}).
\end{equation}

Let \(\lambda_{n,k}^{(q)}\) be the Lagrange weight for evaluating a polynomial
at zero from its values at \(1,q,\ldots,q^n\).  The standard closed form is
\begin{equation}\label{eq:lagrange-weight}
 \lambda_{n,k}^{(q)}
 =\frac{(-1)^{n-k}q^{\binom{n-k+1}{2}}}{(q;q)_n}
 \qbinom{n}{k}{q}.
\end{equation}

\begin{theorem}[Reciprocal row reversal; \Derived]\label{thm:lagrange-reversal}
For every admissible \(q\ne0\),
\begin{equation}\label{eq:lagrange-reversal}
 \lambda_{n,k}^{(q)}=\lambda_{n,n-k}^{(q^{-1})}.
\end{equation}
Hence the rows for \(q=2,-2,4,-4\) are exact reversals of the rows for
\(q=1/2,-1/2,1/4,-1/4\), respectively.
\end{theorem}

\begin{proof}
The geometric node list for \(q^{-1}\), read backwards, is
\(q^{-n},q^{-(n-1)},\ldots,1\), which becomes the \(q\)-node list after a common
nonzero scaling by \(q^n\).  Lagrange weights for evaluation at zero are
unchanged by common scaling and are reversed with the node order.  Alternatively,
substitution of \eqref{eq:qbinom-reciprocity} into
\eqref{eq:lagrange-weight} gives the same identity algebraically.
\end{proof}

\subsection{Finite normalized geometric laws}
The infinite probability series \eqref{eq:Xq-def} does not converge for
\(|q|>1\).  There is nevertheless a canonical finite normalization:
\begin{equation}\label{eq:finite-law}
 X_q^{[m]}=\alpha_{m,q}\sum_{j=0}^{m-1}q^jU_j,
 \qquad
 \alpha_{m,q}=\frac{1-q}{1-q^m},
\end{equation}
provided \(q\ne1\) and \(q^m\ne1\).  The coefficients sum to one, although
negative \(q\) produces signed affine coefficients.

\begin{theorem}[Exact reciprocal-radix law; \Derived]\label{thm:finite-reciprocal}
For every real admissible \(q\ne0\),
\begin{equation}\label{eq:finite-reciprocal}
 X_q^{[m]}\stackrel d=X_{q^{-1}}^{[m]}.
\end{equation}
If \(|q|>1\), then
\begin{equation}\label{eq:outside-limit}
 X_q^{[m]}\Longrightarrow X_{q^{-1}}
 \qquad(m\to\infty).
\end{equation}
Thus the exterior parameters \(2,-2,4,-4\), interpreted through normalized
finite sums, converge to the interior laws at \(1/2,-1/2,1/4,-1/4\).
\end{theorem}

\begin{proof}
A direct calculation shows
\[
 \alpha_{m,q}q^{m-1-j}=\alpha_{m,q^{-1}}q^{-j}.
\]
Reversing the iid list \((U_0,\ldots,U_{m-1})\) therefore turns
\(X_q^{[m]}\) into \(X_{q^{-1}}^{[m]}\) without changing its law.  If
\(|q|>1\), then \(|q^{-1}|<1\), and the finite truncations on the reciprocal
side converge almost surely and in every \(L^p\) to \(X_{q^{-1}}\).
\end{proof}

This theorem is the probability counterpart of Pochhammer reciprocity and
Lagrange row reversal.  It also prevents a common conceptual mistake: values
\(|q|>1\) should not be inserted naively into the infinite series, but they do
possess exact finite and limiting meanings after normalization and digit
reversal.

\section{A Thue--Morse desingularization theorem}

\subsection{Signed subset exponential transforms}
Let \(b_j(n)\in\{0,1\}\) be the binary digits of \(n<2^m\), let
\(s_2(n)=\sum_jb_j(n)\), and put \(\epsTM_n=(-1)^{s_2(n)}\).  For a real radix
\(b\), define the signed digit evaluation
\begin{equation}\label{eq:Ab}
 A_b(n)=\sum_{j=0}^{m-1}b^jb_j(n).
\end{equation}
The finite exponential transform is
\begin{equation}\label{eq:Tmb}
 T_{m,b}(s)=\prod_{j=0}^{m-1}(1-\e^{b^js})
 =\sum_{n=0}^{2^m-1}\epsTM_n\e^{sA_b(n)}.
\end{equation}
For \(b=2\), \(A_2(n)=n\), so this is the classical Thue--Morse exponential
polynomial.  For \(b=-2\), it evaluates the same binary digit subsets in
negabinary place values; \(b=\pm4\) produces sparse base-four and negative
base-four evaluations.

The product in \eqref{eq:Tmb} has a zero of exact order \(m\) at \(s=0\), the
analytic form of Prouhet cancellation.  Dividing out that zero produces not
merely a convenient regularization, but a probability transform.

\begin{theorem}[Thue--Morse desingularization; \Derived]\label{thm:desingularization}
Let \(b\in\R\setminus\{0,1\}\), assume \(b^m\ne1\), and set
\(\alpha_{m,b}=(1-b)/(1-b^m)\).  Then
\begin{equation}\label{eq:desingularization}
 \boxed{
 \E\e^{t(X_b^{[m]}-1/2)}
 =\frac{(-1)^m\e^{-t/2}T_{m,b}(\alpha_{m,b}t)}
 {\alpha_{m,b}^m b^{\binom m2}t^m}.}
\end{equation}
The right side has a removable singularity at \(t=0\), with value one.  For
\(|b|>1\),
\begin{equation}\label{eq:desingularization-limit}
 \E\e^{t(X_{1/b}-1/2)}
 =\lim_{m\to\infty}
 \frac{(-1)^m\e^{-t/2}T_{m,b}(\alpha_{m,b}t)}
 {\alpha_{m,b}^m b^{\binom m2}t^m},
\end{equation}
locally uniformly for complex \(t\).
\end{theorem}

\begin{proof}
The centered moment generating function of \eqref{eq:finite-law} is
\begin{align*}
 \E\e^{t(X_b^{[m]}-1/2)}
 &=\e^{-t/2}\prod_{j=0}^{m-1}
 \frac{\e^{\alpha_{m,b}b^jt}-1}{\alpha_{m,b}b^jt}\\
 &=\frac{(-1)^m\e^{-t/2}}
 {\alpha_{m,b}^mb^{\sum_{j=0}^{m-1}j}t^m}
 \prod_{j=0}^{m-1}(1-\e^{\alpha_{m,b}b^jt}),
\end{align*}
which is \eqref{eq:desingularization}.  The cancellation of the order-\(m\)
zero is explicit in the product.  For \(|b|>1\),
\cref{thm:finite-reciprocal} identifies the finite law with the truncation at
\(1/b\); bounded support on compact \(t\)-sets gives locally uniform convergence
of moment generating functions.
\end{proof}

\subsection{The four requested reciprocal-radix limits}
For \(b=2\), \(\alpha_{m,2}=1/(2^m-1)\) and \(A_2(n)=n\), so
\begin{equation}\label{eq:fabius-thue-limit}
 \psi_{1/2}(-\ii t)
 =\E\e^{t(X_{1/2}-1/2)}
 =\lim_{m\to\infty}
 \frac{(-1)^m\e^{-t/2}}
 {\alpha_{m,2}^m2^{\binom m2}t^m}
 \sum_{n=0}^{2^m-1}\epsTM_n\e^{\alpha_{m,2}tn}.
\end{equation}
Replacing \(t\) by \(\ii t\) gives a characteristic-function formula.  This
is a direct limiting representation of the Fabius/Rvachev Fourier product by
renormalized Thue--Morse exponential sums.

For \(b=-2\), the same theorem converges to \(X_{-1/2}\), which is the affine
stretch of the Rvachev law in \eqref{eq:minus-half-up}.  For \(b=4\) and
\(b=-4\), it converges to the quarter-base laws in
\eqref{eq:poch-quarter}--\eqref{eq:poch-minus-quarter}.  Therefore the eight
values named in the prompt are organized into four reciprocal pairs:
\begin{center}
\begin{tabularx}{0.96\textwidth}{@{}ccc>{\raggedright\arraybackslash}X@{}}
\toprule
exterior radix & reciprocal parameter & Pochhammer base & normalized shape \\
\midrule
\(2\) & \(1/2\) & \(1/4\) & Rvachev \(\Up\) \\
\(-2\) & \(-1/2\) & \(1/4\) & same as \(\Up\), affinely stretched \\
\(4\) & \(1/4\) & \(1/16\) & quarter-base plateau profile \\
\(-4\) & \(-1/4\) & \(1/16\) & same quarter-base profile, affinely stretched \\
\bottomrule
\end{tabularx}
\end{center}

\subsection{Why Bernoulli numbers occur on both sides}
The normalized factor
\[
 \frac{\e^x-1}{x}
\]
is simultaneously the moment generating function of a uniform digit and the
factor obtained after desingularizing \(1-\e^x\).  Hence the same Bernoulli
series \eqref{eq:uniform-cgf} governs both geometric-uniform cumulants and the
Taylor expansion of Thue--Morse products.  More quantitatively,
\begin{equation}\label{eq:qinteger-limit}
 \alpha_{m,b}^{2\ell}[m]_{b^{2\ell}}
 \longrightarrow
 \frac{(1-b^{-1})^{2\ell}}{1-b^{-2\ell}}
 \qquad(|b|>1),
\end{equation}
which is exactly the geometric cumulant factor at \(q=1/b\).  Thus the
Bernoulli denominators \(1-q^{2\ell}\) are reciprocal-radix limits of finite
Gaussian \(q\)-integers.

\section{A digit-position Thue--Morse/\texorpdfstring{\(q\)}{q}-Pochhammer master product}

\subsection{The two-variable product}
Define the binary position weight
\begin{equation}\label{eq:pi2}
 \pi_2(n)=\sum_{j\ge0}j\,b_j(n).
\end{equation}
For \(m\ge0\), introduce
\begin{equation}\label{eq:Mmaster}
 \mathcal M_m(u,z;q)=\prod_{j=0}^{m-1}(1-uq^jz^{2^j}).
\end{equation}

\begin{theorem}[Digit-position master identity; classical ingredients,
new-to-corpus synthesis]\label{thm:digit-master}
One has the exact expansion
\begin{equation}\label{eq:digit-master}
 \mathcal M_m(u,z;q)
 =\sum_{n=0}^{2^m-1}(-u)^{s_2(n)}q^{\pi_2(n)}z^n.
\end{equation}
Consequently,
\begin{equation}\label{eq:poch-digital}
 (u;q)_m=\sum_{n<2^m}(-u)^{s_2(n)}q^{\pi_2(n)},
\end{equation}
and, for each \(0\le k\le m\),
\begin{equation}\label{eq:fixed-weight}
 \sum_{\substack{n<2^m\\s_2(n)=k}}q^{\pi_2(n)}
 =q^{\binom k2}\qbinom{m}{k}{q}.
\end{equation}
\end{theorem}

\begin{proof}
Choosing the second term from the \(j\)-th factor of
\eqref{eq:Mmaster} records the bit \(b_j(n)=1\).  The product of chosen
contributions is
\((-u)^{s_2(n)}q^{\sum jb_j(n)}z^{\sum2^jb_j(n)}\), proving
\eqref{eq:digit-master}.  Set \(z=1\) for \eqref{eq:poch-digital}.  Grouping
subsets by cardinality and using the finite \(q\)-binomial theorem gives
\eqref{eq:fixed-weight}; this fixed-weight identity is also a specialization of
the \(q\)-digital binomial theorem \cite{mansour-nguyen}.
\end{proof}

The ordinary Thue--Morse product is obtained by setting \(u=1,q=1\).  The
ordinary Hamming-weight enumerator is obtained by setting \(q=1\).  The
\(q\)-Pochhammer symbol is obtained by suppressing the integer coordinate
\(z\).  Thus \eqref{eq:digit-master} simultaneously tracks the integer encoded
by a subset, its Hamming weight, and the sum of its occupied bit positions.

\subsection{A nondegenerate signed specialization}
The direct choice \(u=1,z=1\) includes the zero factor at \(j=0\).  Shift the
weight by setting \(u=q\) and define
\begin{equation}\label{eq:Tmaster}
 \mathcal T_m(z;q)=\prod_{j=0}^{m-1}(1-q^{j+1}z^{2^j}).
\end{equation}
Let
\[
 \rho_2(n)=\sum_{j\ge0}(j+1)b_j(n)=\pi_2(n)+s_2(n).
\]
Then
\begin{equation}\label{eq:Tmaster-expansion}
 \mathcal T_m(z;q)=
 \sum_{n<2^m}\epsTM_nq^{\rho_2(n)}z^n,
 \qquad
 \mathcal T_m(1;q)=(q;q)_m.
\end{equation}
At fixed Hamming weight,
\begin{equation}\label{eq:T-fixed}
 \sum_{\substack{n<2^m\\s_2(n)=k}}q^{\rho_2(n)}
 =q^{k(k+1)/2}\qbinom{m}{k}{q}.
\end{equation}

\subsection{Negative \texorpdfstring{\(q\)}{q} creates a position-parity companion}
Let
\begin{equation}\label{eq:chi}
 \chi(n)=(-1)^{\pi_2(n)}.
\end{equation}
For \(q=-r\), \(r>0\), the coefficient in \eqref{eq:Tmaster-expansion}
becomes
\begin{equation}\label{eq:negative-position-parity}
 \epsTM_n(-r)^{\rho_2(n)}
 =(-1)^{s_2(n)+\pi_2(n)+s_2(n)}r^{\rho_2(n)}
 =\chi(n)r^{\rho_2(n)}.
\end{equation}
Thus negative \(q\) transfers the usual Thue--Morse Hamming parity into parity
of the occupied bit positions.

\begin{proposition}[Four-state automatic companion; \Derived]
The pair \((\epsTM_n,\chi(n))\) is \(2\)-automatic.  Its recurrences are
\begin{align}
 \epsTM_{2n}&=\epsTM_n,&\epsTM_{2n+1}&=-\epsTM_n,\label{eq:eps-rec}\\
 \chi(2n)&=\chi(n)\epsTM_n,&
 \chi(2n+1)&=\chi(n)\epsTM_n.\label{eq:chi-rec}
\end{align}
\end{proposition}

\begin{proof}
Shifting the binary digits of \(n\) one place gives
\(\pi_2(2n)=\pi_2(n)+s_2(n)\); adding the new least significant bit to form
\(2n+1\) does not change \(\pi_2\).  The recurrences follow.  The pair takes
only four sign states, and each input bit updates the state deterministically.
\end{proof}

\subsection{Mahler equation and reciprocal symmetry}
The infinite product, when convergent as a formal or analytic object, satisfies
\begin{equation}\label{eq:master-mahler}
 \mathcal M_\infty(u,z;q)
 =(1-uz)\mathcal M_\infty(uq,z^2;q).
\end{equation}
The finite product has the exact reciprocity
\begin{equation}\label{eq:master-reciprocity}
 \mathcal M_m(u,z;q)
 =(-u)^mq^{\binom m2}z^{2^m-1}
 \mathcal M_m(u^{-1},z^{-1};q^{-1}).
\end{equation}
For \(\mathcal T_m\), this becomes
\begin{equation}\label{eq:T-reciprocity}
 \mathcal T_m(z;q)=(-1)^mq^{m(m+1)/2}z^{2^m-1}
 \mathcal T_m(z^{-1};q^{-1}).
\end{equation}
These identities give a digit-polynomial counterpart to the reciprocal
probability law and reciprocal Lagrange rows.

\section{General-radix Bell--Bernoulli formulas}

\subsection{Exact exponential normal form}
Return to \(T_{m,b}\) in \eqref{eq:Tmb}.  Put
\([m]_b=1+b+\cdots+b^{m-1}\).  The uniform cumulant expansion gives an exact
formal normal form.

\begin{theorem}[General-radix Bell--Bernoulli formula; \Derived extension]
For every nonzero \(b\),
\begin{equation}\label{eq:T-normal-form}
 T_{m,b}(t)=(-1)^mb^{\binom m2}t^m
 \exp\left(\gamma_1t+
 \sum_{\ell\ge1}\gamma_{2\ell}\frac{t^{2\ell}}{(2\ell)!}\right),
\end{equation}
where
\begin{equation}\label{eq:gamma}
 \gamma_1=\frac12[m]_b,
 \qquad
 \gamma_{2\ell}=\frac{B_{2\ell}}{2\ell}[m]_{b^{2\ell}},
 \qquad
 \gamma_{2\ell+1}=0\quad(\ell\ge1).
\end{equation}
Consequently, for every \(r\ge0\),
\begin{equation}\label{eq:radix-moment-sum}
 \sum_{S\subseteq\{0,\ldots,m-1\}}(-1)^{|S|}
 \left(\sum_{j\in S}b^j\right)^{m+r}
 =(-1)^mb^{\binom m2}\frac{(m+r)!}{r!}
 \Bell_r(\gamma_1,\ldots,\gamma_r).
\end{equation}
All corresponding signed power sums of degree below \(m\) vanish.
\end{theorem}

\begin{proof}
Factor
\[
 1-\e^x=-x\frac{\e^x-1}{x}
\]
in each term of \(T_{m,b}\), and use \eqref{eq:uniform-cgf}.  Summing
\(b^{nj}\) over \(j\) gives \([m]_{b^n}\), yielding
\eqref{eq:T-normal-form}.  Expanding the exponential with complete Bell
polynomials and comparing the coefficient of \(t^{m+r}\) with the subset-sum
expansion \eqref{eq:Tmb} gives \eqref{eq:radix-moment-sum}.  The factor \(t^m\)
proves the lower-degree cancellations.
\end{proof}

At \(b=2\), this recovers the dyadic Bell--Bernoulli formulas already developed
in the repository.  The same theorem supplies explicit versions at
\(b=-2,4,-4\), and indeed at arbitrary real or complex radix, without a new
calculation.

\subsection{Centered positive/negative-radix conjugacy}
\begin{theorem}[Centered radix sign conjugacy; \Derived]\label{thm:radix-sign}
For \(a>0\),
\begin{equation}\label{eq:radix-sign}
 \e^{-\frac12[m]_{-a}t}T_{m,-a}(t)
 =(-1)^{\binom m2}
 \e^{-\frac12[m]_at}T_{m,a}(t).
\end{equation}
\end{theorem}

\begin{proof}
Since \([m]_b=\sum_{j=0}^{m-1}b^j\), the centered product factors as
\[
 \e^{-\frac12[m]_bt}T_{m,b}(t)
 =\prod_{j=0}^{m-1}\e^{-b^jt/2}(1-\e^{b^jt})
 =\prod_{j=0}^{m-1}(-2\sinh(b^jt/2)).
\]
Replacing \(b=a\) by \(-a\) multiplies the \(j\)-th hyperbolic sine by
\((-1)^j\).  Their product is
\((-1)^{\sum_{j=0}^{m-1}j}=(-1)^{\binom m2}\).
\end{proof}

This is the discrete signed-transform analogue of
\cref{thm:sign-conjugacy}.  In the continuous probability law, each centered
uniform Fourier factor is a sinc, which is even, so coordinate signs disappear
without a residual factor.  In the centered Thue--Morse transform, each factor
is a hyperbolic sine, which is odd, so the coordinate signs accumulate into the
global parity \((-1)^{\binom m2}\).

\section{A unified coordinate-sign and scale-reversal principle}
The principal identities can now be arranged into one diagram:
\[
\begin{array}{ccc}
\text{expanding radix }b & \xleftrightarrow{\text{reverse digits}} &
\text{contracting ratio }q=b^{-1}\\[0.3em]
T_{m,b}\text{ with an order-}m\text{ zero}
&\xrightarrow{\text{desingularize}}&
\text{finite uniform convolution }X_b^{[m]}\\[0.3em]
q\text{-Pochhammer reciprocity}
&\longleftrightarrow&
\text{geometric Lagrange row reversal}\\[0.3em]
\text{odd }\sinh\text{ sign accumulation}
&\longleftrightarrow&
\text{even }\sinc\text{ sign invisibility}.
\end{array}
\]
Three operations repeatedly reappear:
\begin{enumerate}
\item \textbf{Centering} converts uniform moment factors into even sinc or
  hyperbolic-sinc factors.
\item \textbf{Reversing positions} converts exterior scales \(|b|>1\) into
  interior ratios \(|q|<1\).
\item \textbf{Splitting positions} converts \(q\) into \(q^m\), simultaneously
  dissection of random digits, Gaussian \(q\)-integers, and Pochhammer products.
\end{enumerate}
This principle explains why \(q^2\), rather than \(q\), is the natural base in
the centered Fourier Pochhammer factorization: centering makes the factor even,
and evenness forgets the sign of every geometric scale.

\section{Lambert--\texorpdfstring{\(W\)}{W} endpoint asymptotics beyond \texorpdfstring{\(q=1/2\)}{q=1/2}}

\subsection{The positive-parameter Laplace cocycle}
For \(0<r<1\), the Laplace transform is
\begin{equation}\label{eq:laplace-product}
 L_r(s)=\E\e^{-sX_r}
 =\prod_{n=0}^{\infty}
 \frac{1-\e^{-(1-r)r^ns}}{(1-r)r^ns},
 \qquad s>0.
\end{equation}
It satisfies the exact functional equation
\begin{equation}\label{eq:laplace-functional}
 L_r(s)=\frac{1-\e^{-(1-r)s}}{(1-r)s}L_r(rs).
\end{equation}
Set
\[
 \lambda=-\log r,
 \qquad
 y=\log((1-r)s),
 \qquad
 G_r(y)=\log L_r\!\left(\frac{\e^y}{1-r}\right).
\]
Then
\begin{equation}\label{eq:additive-cocycle}
 G_r(y)-G_r(y-\lambda)
 =h(y),
 \qquad
 h(y)=\log\frac{1-\e^{-\e^y}}{\e^y}.
\end{equation}
As \(y\to+\infty\), \(h(y)=-y+\bigO(\e^{-\e^y})\).  The quadratic function
\[
 Q_r(y)=-\frac{y^2}{2\lambda}-\frac y2
\]
satisfies \(Q_r(y)-Q_r(y-\lambda)=-y\).  This forces the leading log-square
behavior behind the known half-base Lambert analysis in the repository.

\begin{conjecture}[Periodic cocycle normal form; \Conjectural]\label{conj:periodic-cocycle}
For every \(0<r<1\), there is a real-analytic one-periodic function \(P_r\) and
a constant \(C_r\) such that
\begin{equation}\label{eq:periodic-cocycle-conj}
 G_r(y)=-\frac{y^2}{2\lambda}-\frac y2+C_r
 +P_r\!\left(\frac y\lambda\right)+\smallo(1)
 \qquad(y\to+\infty),
\end{equation}
with a full asymptotic expansion obtained by solving the additive cohomological
equation \eqref{eq:additive-cocycle}.  The Fourier coefficients of \(P_r\)
should be expressible through the Mellin transform of
\(h(y)+y\mathbf 1_{y>0}\).
\end{conjecture}

The periodic term is not cosmetic.  It records the fractional position of the
saddle relative to the geometric lattice \(r^n\), exactly as dyadic endpoint
asymptotics retain a lower-Lambert phase in the current frontier corpus.

\subsection{The Lambert saddle}
Ignoring the bounded periodic derivative at first order, the saddle equation
for Laplace inversion has the form
\begin{equation}\label{eq:saddle-leading}
 x\sim\frac{(1-r)y}{\lambda\e^y}.
\end{equation}
Therefore
\begin{equation}\label{eq:lambert-leading}
 y\sim-W_{-1}\!\left(-\frac{\lambda}{1-r}x\right),
 \qquad x\downarrow0,
\end{equation}
where \(W_{-1}\) is the lower real branch of Lambert's function.

\begin{conjecture}[All-order general-\(q\) endpoint transseries;
\Conjectural]\label{conj:lambert-allorder}
For each \(0<r<1\), the density, CDF, and inverse CDF near the left endpoint
admit all-order transseries in inverse powers of
\[
 w_r(x)=-W_{-1}\!\left(-\frac{-\log r}{1-r}x\right),
\]
whose coefficients are periodic functions of \(w_r(x)/(-\log r)\).  The
coefficient algebra should be generated by Bell polynomials in the derivatives
of the cocycle correction \(P_r\), and should reduce at \(r=1/2\) to the
existing dyadic Lambert-phase machinery.
\end{conjecture}

\subsection{Negative endpoints require no new saddle calculation}
Let \(a_{-r}=-r/(1-r)\) be the left endpoint of \(X_{-r}\).  From
\eqref{eq:cdf-transport},
\begin{equation}\label{eq:negative-endpoint-transfer}
 F_{-r}(a_{-r}+\delta)=F_r(\delta/c_r).
\end{equation}
Thus every proved or conjectured positive-\(r\) endpoint expansion transfers
exactly to the negative parameter.  In particular, the relevant Lambert
variable is simply \(w_r(\delta/c_r)\).  The inverse endpoint expansion obeys
the affine formula \eqref{eq:quantile-transport}.  This is substantially
stronger than a parallel asymptotic derivation: there is no independent
negative-\(q\) transseries to compute.

\section{Numerical experiments}

The accompanying script \texttt{q\_fabius\_experiments.py} is deterministic and
contains detailed comments.  It uses exact rational arithmetic for finite
\(q\)-binomial, Lagrange, and digit-product identities; 80-digit arithmetic for
analytic product comparisons; and a padded \(2^{16}\)-point FFT for density
inversion.  The complete output is in \texttt{numerical\_checks.csv}.

\subsection{Selected verification results}
\begin{center}
\begin{tabularx}{0.96\textwidth}{@{}>{\raggedright\arraybackslash}p{0.34\textwidth}
>{\raggedright\arraybackslash}p{0.26\textwidth}
>{\raggedleft\arraybackslash}p{0.16\textwidth}X@{}}
\toprule
identity & parameters & residual & metric \\
\midrule
sign conjugacy & \(r=1/2\) & \(0\) & max characteristic error \\
\(q\)-power decimation & \(q=1/2,m=3\) & \(6.7\cdot10^{-16}\) & max characteristic error \\
sinc--Pochhammer logarithm & \(q=-1/2,t=1.1\) & \(3.7\cdot10^{-72}\) & 80-digit log error \\
Lagrange reciprocal reversal & all \(q=\pm1/2,\pm2,\pm1/4,\pm4\), \(n=3,5,7\) & \(0\) & exact rational mismatch count \\
digit-position product & \(m=8,u=2/3,q=-1/2\) & \(0\) & exact coefficient mismatches \\
standardized sign collapse & \(r=3/4\) & \(4.8\cdot10^{-16}\) & max FFT density error \\
quarter-base plateau & \(q=1/4\) & \(6.7\cdot10^{-16}\) & max interior height error \\
\bottomrule
\end{tabularx}
\end{center}

\begin{figure}[H]
\centering
\includegraphics[width=0.96\textwidth]{physical_densities.pdf}
\caption{Physical densities for \(q=\pm1/2,\pm1/4\).  The negative curves are
exact affine stretches of their positive partners.  The quarter-base profiles
have positive-length plateaus; the half-base plateau is critical and collapses
to the center.}
\label{fig:physical-densities}
\end{figure}

\begin{figure}[H]
\centering
\includegraphics[width=0.96\textwidth]{normalized_profiles.pdf}
\caption{Support-normalized profiles \(\mathcal U_q\) for positive \(q\).
The \(q=1/4\) plateau, the critical Rvachev profile at \(q=1/2\), and the
increasingly Gaussian central concentration as \(q\to1^-\) are visible.
The profiles for negative \(q\) are identical after support normalization.}
\label{fig:normalized-profiles}
\end{figure}

\subsection{Interpretation and limitations}
The numerical work verifies formulas already proved analytically; it is not
used as a substitute for proof.  FFT inversion is most accurate away from exact
support endpoints, where the true density is infinitely flat but floating-point
cancellation is severe.  The exact plateau and sign-collapse checks are
performed on interior grids.  The script retains tiny negative values of order
machine precision rather than clipping them, so the diagnostic CSV exposes
roundoff instead of hiding it.

\section{Further conjectures and research directions}

\subsection{A decimation renormalization semigroup}
The maps \(q\mapsto q^m\) and weights \(a_{m,j}(q)\) satisfy a refinement
associativity obtained by splitting indices first modulo \(m\) and then modulo
\(n\), or directly modulo \(mn\).  This suggests a semigroup of exact
convolution factorizations acting on densities, cumulants, spectral zeta
functions, and Pochhammer products.

\begin{problem}[Optimal decimation algorithms]
Develop stable exact and numerical algorithms that choose \(m\) adaptively to
represent a difficult profile at \(q\) by simpler profiles at \(q^m\).  For
\(|q|\) near one, repeated powering rapidly enters the plateau regime
\(|q^m|<1/2\), where the component density has an exact central constant region.
Quantify conditioning, Fourier truncation error, and optimal convolution order.
\end{problem}

\subsection{Two-base Mahler systems}
The master product \(\mathcal M_m(u,z;q)\) couples a geometric position weight
\(q^j\) to a lacunary exponent \(2^j\).  Replacing \(2\) by an integer radix
\(b\ge2\) gives
\[
 \mathcal M_\infty^{(b)}(u,z;q)
 =\prod_{j\ge0}(1-uq^jz^{b^j}),
 \qquad
 \mathcal M_\infty^{(b)}(u,z;q)
 =(1-uz)\mathcal M_\infty^{(b)}(uq,z^b;q).
\]

\begin{conjecture}[Joint natural boundaries; \Conjectural]
For generic algebraic \((q,z)\) with \(0<|q|<1\) and \(|z|=1\), the two-base
product has a natural boundary in the \(z\)-plane unless the parameters lie on
an explicitly classifiable cyclotomic or telescoping locus.  The exceptional
loci should be controlled by multiplicative dependence between \(q\) and
powers of \(z\).
\end{conjecture}

A useful first step is to derive a complete matrix Mahler system for the
four-state pair \((\epsTM_n,\chi(n))\), then study its Dirichlet and Mellin
transforms in the spirit of \cite{toth}.

\subsection{Strict shape and total positivity}
Log-concavity is only the first layer of possible shape control.

\begin{problem}[P\'olya-frequency structure]
Determine for which \(q\) the density \(f_q\), or its sampled Toeplitz kernels,
has finite or infinite total positivity.  Finite convolutions of interval
indicators are box splines, but the geometric limit may retain stronger
variation-diminishing properties.  Establishing a suitable strict total
positivity theorem could resolve strict log-concavity and uniqueness of the
mode for \(|q|>1/2\).
\end{problem}

\subsection{Arithmetic of quarter-base and signed dyadic values}
Arias de Reyna and Haugland establish rich arithmetic at dyadic points for the
ordinary Fabius function \cite{arias-arithmetic,haugland}.  Decimation suggests
that quarter-base values are not merely analogues but convolutional building
blocks of the dyadic law.

\begin{problem}[Arithmetic transfer through decimation]
Investigate whether denominators, valuations, and integrality patterns for
\(F_{1/2}\) can be decomposed into bilinear or multinomial expressions in
values and moments of \(F_{1/4}\).  The same question at
\(q=2^{-m}\) may reveal a hierarchy of arithmetic invariants compatible with
Gaussian \(q\)-binomial coefficients and the weight vector in
\eqref{eq:half-to-2m}.
\end{problem}

For negative parameters, all value questions reduce affinely to positive
parameters.  For example, rationality or transcendence of
\(F_{-1/2}(x)\) is exactly the corresponding question for
\(F((x+1)/3)\).  This should be used to avoid duplicating arithmetic work.

\subsection{Complex parameters and roots of unity}
The probability interpretation is real, but the finite algebraic identities
remain meaningful for complex \(q\).  Near roots of unity, denominators in
Gaussian interpolation and finite normalizations become singular, while
Pochhammer products exhibit cyclotomic cancellation.

\begin{problem}[Root-of-unity regularization]
Construct canonical renormalized limits of the finite geometric laws,
Lagrange rows, and digit-position products as \(q\) approaches a root of unity.
Classify which singularities are removable after centering, residue-class
splitting, or grouping full cyclotomic orbits.
\end{problem}

\subsection{General-\texorpdfstring{\(q\)}{q} endpoint analysis}
The most promising analytic program is to turn
\cref{conj:periodic-cocycle,conj:lambert-allorder} into theorems.  A practical
route is:
\begin{enumerate}
\item solve \eqref{eq:additive-cocycle} by a convergent Mellin--Fourier
  decomposition after subtracting the quadratic model;
\item obtain uniform bounds for derivatives of the periodic correction;
\item perform the saddle analysis with the periodic phase retained from the
  start, rather than inserted as a late perturbation;
\item encode the all-order coefficient recursion with complete and partial Bell
  polynomials;
\item transfer the result to negative \(q\) by
  \eqref{eq:negative-endpoint-transfer} and to inverse CDFs by affine quantile
  transport.
\end{enumerate}

\section{Formalization roadmap}
The exact results in this report are unusually suitable for formalization
because their proofs reduce to reindexing summable series, symmetry of the
uniform law, finite products, and standard characteristic-function facts.
A staged Lean plan is as follows.

\subsection{Stage A: sign conjugacy}
Formalize the centered uniform symmetry and the deterministic sign-twist of an
iid sequence.  Prove equality in distribution of the two weighted series, then
derive pushforward formulas for CDFs, densities, support, and quantiles.  This
would close the negative-parameter density gap identified in the current
frontier register.

\subsection{Stage B: residue-class decimation}
Build an equivalence between \(\N\) and \(\{0,\ldots,m-1\}\times\N\), prove
absolute summability under reindexing, and establish independence of the
subsequences.  Characteristic and cumulant corollaries should then be short.

\subsection{Stage C: Fourier decay and smoothness}
Formalize \(|\sinc x|\le\min(1,1/|x|)\), the logarithmic choice of the cutoff,
and integrability of all polynomial multiples.  Reuse the existing Fourier
inversion infrastructure to obtain \(C^\infty\) densities for negative \(q\)
and endpoint flatness.

\subsection{Stage D: finite algebra}
The Pochhammer reciprocity, Gaussian reciprocity, Lagrange reversal, digital
master product, and Bell--Bernoulli finite identities are polynomial or finite
product statements.  They can be proved before any analytic infrastructure.
The exact rational checks in the accompanying script supply regression data,
not proof certificates.

\subsection{Stage E: Thue--Morse desingularization}
Define \(T_{m,b}\) both as a product and a Boolean-cube sum, prove the order-
\(m\) zero, and identify the regularized quotient with the finite centered
uniform transform.  The reciprocal finite-law theorem then supplies the limit
at \(|b|>1\).

\section{Conclusion}
The requested values \(\pm1/2,\pm1/4,\pm2,\pm4\) do not form eight unrelated
cases.  They sit in a compact web generated by three exact symmetries:
\[
 q\longmapsto-q,
 \qquad
 q\longmapsto q^m,
 \qquad
 q\longmapsto q^{-1}.
\]
The first is affine sign conjugacy of centered uniform digits.  The second is
residue-class decimation and simultaneously Pochhammer dissection.  The third
is finite digit reversal, Gaussian/Pochhammer reciprocity, and Lagrange row
reversal.  Their combination yields the following central picture.

\begin{itemize}
\item The negative half-base and quarter-base densities are exact affine
  stretches of the positive profiles; after standardization or support
  normalization, the sign of \(q\) disappears.
\item The Fabius/Rvachev law at \(q=1/2\) is an explicit convolution of two
  quarter-base laws, and more generally a barycenter of \(q=2^{-m}\) laws.
\item The centered Fourier transform is a product of \(q^2\)-Pochhammer
  symbols indexed by the ordinary sine zeros, making the spectral zeta and
  Bernoulli cumulants immediate.
\item Normalized finite laws at \(q=2,-2,4,-4\) are exactly reciprocal to the
  interior finite laws and converge to the requested Fabius-type limits.
\item Desingularizing the order-\(m\) zero of a finite Thue--Morse exponential
  polynomial produces exactly the centered transform of a finite
  geometric-uniform convolution.  In the reciprocal-radix limit, this becomes
  the Fabius/Rvachev transform.
\item Negative \(q\) in the digit-position Pochhammer product replaces the
  Thue--Morse Hamming-parity sign by a four-state position-parity companion.
\item The whole real family is log-concave; \(|q|\le1/2\) is an exact plateau
  phase, while \(|q|\to1\) gives a standardized Gaussian limit.
\end{itemize}

The strongest next step is a rigorous all-\(q\) Lambert--\(W\) endpoint theory
that retains the geometric periodic phase.  The algebraic and probabilistic
identities developed here substantially reduce that task: only positive
\(0<q<1\) needs independent asymptotic analysis, and every negative case follows
by exact affine transport.

\appendix

\section{Detailed derivation of the logarithmic Pochhammer series}
For \(|t|\) sufficiently small, apply
\(\log(a;Q)_\infty=-\sum_{m\ge1}a^m/(m(1-Q^m))\) to
\eqref{eq:sinc-poch}:
\begin{align}
 \log\psi_q(t)
 &=-\sum_{k\ge1}\sum_{m\ge1}
 \frac1{m(1-q^{2m})}
 \left(\frac{(1-q)^2t^2}{4\pi^2k^2}\right)^m\\
 &=-\sum_{m\ge1}\frac{\zeta(2m)}{m(1-q^{2m})}
 \left(\frac{(1-q)t}{2\pi}\right)^{2m}.
 \label{eq:log-poch-expanded}
\end{align}
Euler's evaluation
\[
 \zeta(2m)=(-1)^{m+1}
 \frac{B_{2m}(2\pi)^{2m}}{2(2m)!}
\]
turns \eqref{eq:log-poch-expanded} into
\[
 \log\psi_q(t)
 =\sum_{m\ge1}(-1)^m
 \frac{B_{2m}}{2m}\frac{(1-q)^{2m}}{1-q^{2m}}
 \frac{t^{2m}}{(2m)!},
\]
which is exactly the centered characteristic cumulant series from
\eqref{eq:cumulants}.  This calculation is the shortest algebraic route from
Pochhammer products to Bernoulli cumulants.

\section{Selected exact \texorpdfstring{\(q\)}{q}-Lagrange rows}
The archive includes \texttt{q\_lagrange\_weights.csv}, containing exact
fractions for \(n=5\) and every \(q\in\{\pm1/2,\pm2,\pm1/4,\pm4\}\).  For
illustration, at \(q=1/2\),
\[
 (\lambda_{5,0},\ldots,\lambda_{5,5})
 =\left(
 -\frac1{9765},\frac2{315},-\frac8{63},
 \frac{64}{63},-\frac{1024}{315},\frac{32768}{9765}
 \right),
\]
and the \(q=2\) row is the exact reversal.  The alternating and quarter-base
rows grow more rapidly, but the reciprocal reversal remains exact over the
rationals.

\section{File inventory and reproducibility}
The ZIP supplement contains:
\begin{itemize}
\item \texttt{fabius\_q\_connections\_report.tex}: this source;
\item \texttt{fabius\_q\_connections\_report.pdf}: the rendered report;
\item \texttt{q\_fabius\_experiments.py}: commented deterministic experiments;
\item \texttt{numerical\_checks.csv}: all residuals and exact mismatch counts;
\item \texttt{q\_lagrange\_weights.csv}: exact reciprocal Lagrange rows;
\item \texttt{standardized\_cumulants.csv}: selected closed-form cumulants;
\item \texttt{physical\_densities.pdf/.png} and
  \texttt{normalized\_profiles.pdf/.png}: figures;
\item \texttt{NUMERICAL\_README.txt}: execution instructions and package
  versions.
\end{itemize}
Run the numerical supplement in its directory with
\begin{verbatim}
python q_fabius_experiments.py
\end{verbatim}
No network access, random seed, or external data file is required.

\begin{thebibliography}{99}

\bibitem{repo-primary}
V.~Reshetnikov,
\emph{Fabius Function and Rvachev Up}, current formal exposition in the
\texttt{ProveIt} repository,
\url{https://github.com/VladimirReshetnikov/ProveIt/tree/main/Analysis/FabiusFunction/docs/Fabius_Function_and_Rvachev_Up},
accessed 28 August 2026.

\bibitem{repo-frontiers}
V.~Reshetnikov,
\emph{Non-Formalized Research Frontiers for the Fabius Function}, consolidated
LaTeX volume,
\url{https://github.com/VladimirReshetnikov/ProveIt/blob/main/Analysis/FabiusFunction/docs/non-formalized-research-frontiers/non-formalized-research-frontiers.tex},
accessed 28 August 2026.

\bibitem{repo-qseries}
V.~Reshetnikov,
\emph{Exponents and \(q\)-Series Frontiers}, draft synthesis,
\url{https://github.com/VladimirReshetnikov/ProveIt/blob/main/Analysis/FabiusFunction/docs/non-formalized-research-frontiers/drafts/exponents-and-q-series/Exponents_and_q_Series_Frontiers/Exponents_and_q_Series_Frontiers.tex},
accessed 28 August 2026.

\bibitem{repo-thuemorse}
V.~Reshetnikov,
\emph{Thue--Morse Atlas and Frontiers}, draft synthesis,
\url{https://github.com/VladimirReshetnikov/ProveIt/blob/main/Analysis/FabiusFunction/docs/non-formalized-research-frontiers/drafts/thue-morse/Thue_Morse_Atlas_and_Frontiers/Thue_Morse_Atlas_and_Frontiers.tex},
accessed 28 August 2026.

\bibitem{repo-archive}
V.~Reshetnikov,
\emph{Fabius-function documentation archive and provenance map},
\url{https://github.com/VladimirReshetnikov/ProveIt/tree/main/Analysis/FabiusFunction/docs/archive},
accessed 28 August 2026.

\bibitem{arias-up}
J.~Arias de Reyna,
\emph{An infinitely differentiable function with compact support: Definition
and properties}, arXiv:1702.05442 (English translation of
Rev. Real Acad. Ciencias Madrid \textbf{76} (1982), 21--38), 2017.

\bibitem{arias-arithmetic}
J.~Arias de Reyna,
\emph{Arithmetic of the Fabius function}, arXiv:1702.06487v3, 2017.

\bibitem{haugland}
J.~K.~Haugland,
\emph{Evaluating the Fabius function}, arXiv:1609.07999v2, 2020.

\bibitem{nguyen-digital}
H.~D.~Nguyen,
\emph{A Digital Binomial Theorem}, arXiv:1412.3181, 2014.

\bibitem{mansour-nguyen}
T.~Mansour and H.~D.~Nguyen,
\emph{A \(q\)-Digital Binomial Theorem}, arXiv:1506.07945, 2015.

\bibitem{aistleitner-hofer-larcher}
C.~Aistleitner, R.~Hofer, and G.~Larcher,
\emph{On parametric Thue--Morse sequences and lacunary trigonometric products},
arXiv:1502.06738v2, 2015.

\bibitem{toth}
L.~T\'oth,
\emph{Linear Combinations of Dirichlet Series Associated with the Thue--Morse
Sequence}, \emph{INTEGERS} \textbf{22} (2022), Article 98;
arXiv:2211.13570.

\bibitem{dlmf}
NIST Digital Library of Mathematical Functions,
\emph{Chapter 17: \(q\)-Hypergeometric and Related Functions}, especially
\S17.2 on \(q\)-Pochhammer symbols and Gaussian polynomials,
\url{https://dlmf.nist.gov/17.2}, accessed 28 August 2026.

\bibitem{allouche-shallit}
J.-P.~Allouche and J.~Shallit,
\emph{The ubiquitous Prouhet--Thue--Morse sequence}, in
C.~Ding, T.~Helleseth, and H.~Niederreiter (eds.),
\emph{Sequences and Their Applications}, Springer, 1999, pp.~1--16.

\end{thebibliography}

\end{document}
